%%=============================================================================
%% Inleiding
%%=============================================================================

\chapter{\IfLanguageName{dutch}{Inleiding}{Introduction}}%
\label{ch:inleiding}

Privileged access verwijst naar toegang tot IT-systemen met verhoogde rechten, waarbij bepaalde gebruikers uitgebreide rechten krijgen om systeemconfiguraties te wijzigen en kritieke infrastructuur te beheren(\cite{koot2024pam}). Dergelijke toegang wordt typisch toegekend aan systeem- en netwerkbeheerders. Omdat deze accounts uitgebreide bevoegdheden hebben, vormen ze een aantrekkelijk doelwit voor cyberaanvallen en vormen ze tegelijk een potentieel risico bij interne fouten of misbruik (\cite{keepersecurity2025risicos}). Het correct beheren en beveiligen van privileged access is daarom een essentieel onderdeel van moderne cybersecurity (\cite{keepersecurity2025risicos}). 


\section{\IfLanguageName{dutch}{Probleemstelling}{Problem Statement}}%
\label{sec:probleemstelling}

Binnen hedendaagse IT-omgevingen maken beheerders dagelijks gebruik van protocollen zoals Secure Shell (SSH) en Remote Desktop Protocol (RDP) om servers en andere systemen te beheren (\cite{rfc4251}, \cite{microsoft_rdp}). Deze vormen van toegang vereisen doorgaans verhoogde rechten en worden vaak gebruikt in kritieke infrastructuren. Wanneer dergelijke toegang onvoldoende wordt gecontroleerd, gelogd of langdurig actief blijft, neemt het risico op beveiligingsincidenten aanzienlijk toe. Dit kan leiden tot ongeautoriseerde wijzigingen, datalekken of zelfs volledige compromitteren van systemen zoals beschreven in figuur \ref{fig:attackchain}.

\begin{figure}[H]

	\centering
	\includegraphics[width=0.8\textwidth]{attackchain.png}
	\caption[Understanding the attack chain (\cite{cyberark_blueprint}).]{\label{fig:attackchain}Understanding the attack chain (\cite{cyberark_blueprint}).}
\end{figure}


De aanleiding voor dit onderzoek ontstond tijdens gesprekken met medewerkers van Ahold Delhaize Group. Binnen de organisatie werd een uitgebreid Privileged Access Management-programma uitgewerkt om privileged toegang binnen een complexe internationale IT-omgeving beter te controleren en te standaardiseren. Daarbij werd niet alleen gefocust op technische beveiliging, maar ook op governance, compliance, auditing, onboarding en gecontroleerde toegang tot kritieke systemen. De operationele en strategische uitdagingen binnen deze omgeving vormden een belangrijke inspiratiebron voor de casestudies binnen deze bachelorproef. 

Privileged Access Management biedt een oplossing door processen die zich richten op het beveiligen, beheren en monitoren van administrator accounts (\cite{garbis2021zero}). Deze oplossingen maken het mogelijk om toegangsrechten tijdelijk toe te kennen, sessies te registreren, wachtwoorden centraal te beheren en afwijkend gedrag sneller te detecteren. Hierdoor kunnen organisaties zowel hun cybersecurity verhogen als beter inzicht krijgen in het gebruik van kritieke toegangsrechten.


\section{\IfLanguageName{dutch}{Onderzoeksvraag}{Research question}}%
\label{sec:onderzoeksvraag}

Op basis van bovenstaande probleemstelling wordt volgende onderzoeksvraag geformuleerd:
Hoe kan een Privileged Access Management tool bijdragen aan de beveiliging van servertoegang via SSH en RDP binnen een bedrijfsnetwerk?
Om deze onderzoeksvraag te beantwoorden, worden volgende deelvragen onderzocht:

\begin{itemize}
	\item Welke beveiligingsrisico’s zijn verbonden aan langdurige en ongelogde beheertoegang tot servers via SSH en RDP?
	\item Welke functionaliteiten bieden Privileged Access Management oplossingen om tijdelijke en gelogde beheertoegang af te dwingen?
	\item In welke mate verschillen PAM oplossingen van elkaar op het vlak van sessielogging en sessierecording?
	\item Hoe ondersteunen PAM oplossingen Just-in-Time toegang en het principe van Zero Standing Privileges?
	\item In welke mate kan een PAM oplossing praktisch worden toegepast om beheertoegang tot servers via SSH en RDP te beveiligen binnen een proof-of-concept?
\end{itemize}

\section{\IfLanguageName{dutch}{Onderzoeksdoelstelling}{Research objective}}%
\label{sec:onderzoeksdoelstelling}

Het doel van deze bachelorproef is om te onderzoeken hoe beheertoegang tot servers beter kan worden beveiligd door het gebruik van PAM tools. Meer specifiek, op beheertoegang tot servers via SSH en RDP binnen on-premise bedrijfsnetwerken. Dit gebeurt enerzijds via een literatuurstudie en marktanalyse van bestaande PAM tools. Hierin worden marktleiders, niche spelers en uitdagers vergeleken.  Anderzijds via een proof of concept waarin twee PAM oplossingen praktisch worden getest en vergeleken. De geteste PAM oplossingen zijn meer gericht naar kleine en middelgrote bedrijven. Op basis daarvan wordt nagegaan welke functionaliteiten relevant zijn voor organisaties en welke oplossingen het meest geschikt zijn voor verschillende IT-omgevingen. Hierbij wordt eveneens onderzocht hoe enterprise-organisaties zoals Ahold Delhaize PAM-capabilities prioriteren op basis van risico, governance, operationele impact en implementatiecomplexiteit.

Het onderzoek wordt als succesvol beschouwd indien:
\begin{itemize}
	\item 	De belangrijkste beveiligingsrisico’s rond SSH- en RDP-beheertoegang duidelijk geïdentificeerd en onderbouwd worden.
	\item	Een objectieve vergelijking van PAM oplossingen wordt opgesteld op basis van vooraf gedefinieerde criteria.
	\item	In de proof of concept aangetoond wordt dat:
	\begin{itemize}	
		\item	Langdurige beheerdersrechten kunnen worden vermeden.
		\item	Just-in-Time toegang effectief kan worden afgedwongen.
		\item	Alle beheerderssessies worden gelogd en, indien mogelijk, opgenomen.
		\item	Directe, onbeheerde toegang tot servers kan worden geblokkeerd.
	\end{itemize}
	\item	De praktische verschillen tussen de geteste PAM oplossingen duidelijk worden geanalyseerd.
	\item	De centrale onderzoeksvraag onderbouwd en concreet wordt beantwoord.
	
	
\end{itemize}

\section{\IfLanguageName{dutch}{Opzet van deze bachelorproef}{Structure of this bachelor thesis}}%
\label{sec:opzet-bachelorproef}

% Het is gebruikelijk aan het einde van de inleiding een overzicht te
% geven van de opbouw van de rest van de tekst. Deze sectie bevat al een aanzet
% die je kan aanvullen/aanpassen in functie van je eigen tekst.

Deze bachelorproef is als volgt opgebouwd. 
In hoofdstuk~\ref{ch:stand-van-zaken}  wordt de stand van zaken besproken, waarbij de relevante literatuur, internationale richtlijnen en marktanalyses rond Privileged Access Management worden geanalyseerd. Dit hoofdstuk vormt de theoretische basis voor het verdere onderzoek.


In hoofdstuk~\ref{ch:methodologie} wordt de methodologie toegelicht. Hierin wordt beschreven hoe het onderzoek werd opgezet, welke fasen werden doorlopen en welke onderzoekstechnieken werden toegepast.

In hoofdstuk~\ref{ch:casestudie} worden twee praktijkgerichte casestudies uitgewerkt. De eerste casestudie behandelt de strategische heroriëntatie van de PAM-omgeving binnen Ahold Delhaize en onderzoekt hoe PAM-capabilities worden geprioriteerd op basis van governance, beveiliging, compliance en operationele efficiëntie. De tweede casestudie focust op de operationele implementatie van CyberArk binnen een storageomgeving en vertaalt deze strategische keuzes naar concrete technische vereisten en dagelijkse beheerprocessen.

Hoofdstuk~\ref{ch:vergelijkende-Studie} bevat de vergelijkende studie van geselecteerde PAM vendors op basis van literatuur en marktanalyses. In dit hoofdstuk worden de functionele verschillen tussen de oplossingen systematisch geanalyseerd.


In hoofdstuk~\ref{ch:proof-of-concept} wordt de proof of concept uitgewerkt, waarin twee PAM oplossingen technisch worden geïmplementeerd en getest binnen een gecontroleerde omgeving.


Hoofdstuk~\ref{ch:vergelijkende-Studie-met-poc} bespreekt de resultaten van de proof of concept en vergelijkt beide oplossingen op basis van vooraf bepaalde criteria.

Tot slot worden in hoofdstuk~\ref{ch:conclusie} de conclusies geformuleerd, wordt de centrale onderzoeksvraag beantwoord en worden aanbevelingen gedaan voor praktijktoepassing en verder onderzoek.
